\chapter{External computed results}
\label{chap:computed-inputs}

This chapter contains the concrete Lin outputs, appendix rows, and finite evidence certificates carried as external computation inputs.  Their schemas were defined earlier.

\section{Lin catalogues and appendix tables}

% SOURCE: aimpaper/112.tex:5-1029, Figure Fig:near126; included at aimpaper/main.tex:2152; local PDF p.43.
\begin{definition}[Near-$126$ figure contract]
  \label{def:near126-figure-contract}
  \uses{thm:appendix-all-rows-encoded}
  \notready
  Figure 9 is a derived visualization over stems $122$--$127$ and filtrations
  $0$--$19$.  Its left panel has 289 dots, 384 black structural edges, and 157
  colored known differentials (87 $d_2$, 44 $d_3$, 17 $d_4$, seven $d_5$,
  and two $d_7$).  Its right panel has 42 dots, 21 black edges, and 12 dashed
  shortest-potential differentials.  Dashed arrows are unresolved candidates,
  and the figure is not accepted as independent evidence because its source
  has no complete legend and uses layout offsets for multiple classes.  Lines
  1031--1919 of \texttt{112.tex} are commented duplicates and are ignored.
\end{definition}

% SOURCE: aimpaper/main.tex:213-240 and 2785-2800; LWXMachine and LWXZenodo bibliography entries.
\begin{proposition}[Lin computation provenance]
  \label{prop:lin-computation-provenance}
  \uses{def:near126-computation-input,def:lin-computation-bundle,def:appendix-evidence-record,def:source-ref}
  \notready
  For a near-$126$ computation input, the spectrum, $E_2$-page, map, initial
  $d_2$, and propagated-output catalogues are tied to version
  \texttt{v126.3.cw49} of Zenodo record \texttt{14875701}.  Each archive file,
  program revision or checksum, table row, and proof identifier is populated
  and mutually consistent.  Until every field is populated, the input remains
  not ready; the interactive plot alone is not sufficient provenance.
\end{proposition}

\begin{proof}
  Cross-check the paper bibliography against the local source inventory, open
  the archived computation files, and record stable paths and hashes.  Verify
  every row key against the artifact that proves or bounds its status.  Reject
  records justified only by visual coordinates or an unversioned generator
  name.
\end{proof}

% SOURCE: aimpaper/main.tex:2785-2792, three manually supplied differentials.
\begin{proposition}[Manual differential inputs]
  \label{prop:appendix-manual-inputs}
  \uses{def:near126-computation-input,prop:lin-computation-provenance,def:external-result,def:external-evidence}
  \notready
  For a near-$126$ computation input $I$, its manual-input ledger separately
  records
  \[
  \begin{aligned}
    d_5(h_0^{24}h_6)&=h_0^2P^6d_0,\\
    d_6(h_0^{55}h_7)&=h_0^2x_{126,60},\\
    d_3(v_2^{16})&=\beta^5g\quad\text{in }tmf.
  \end{aligned}
  \]
  Each is an explicit external result or evidence input with a theorem, page,
  or archived computation locator.  The two image-of-$J$ locators and the
  precise Bruner--Rognes locator are currently provenance obligations.
\end{proposition}

\begin{proof}
  Type-check the three source and target degrees, attach their independent
  source locators, and verify that the archived program declares exactly these
  inputs.  Do not allow a table output derived from an input to serve as the
  source of that same input.
\end{proof}

% SOURCE: LWXMachine Sections 2--3 and Tables of Proofs; Zenodo record
% 14875701, version v126.3.cw49; PROJECT_BOUNDARY.md, Appendix data.
\begin{theorem}[Lin input/output catalogues are complete]
  \label{thm:lin-computation-catalogues-complete}
  \uses{def:near126-computation-input,def:lin-computation-bundle,prop:lin-computation-provenance,prop:appendix-manual-inputs}
  \notready
  For a near-$126$ computation input, every one of the 49 CW spectra, its stored
  $E_2$ generators, relations, and basis through the declared bound, all 180
  maps, every specified initial $d_2$, and every imported propagated
  differential, extension, disproof, or unresolved candidate has exactly one
  typed catalogue record.  Conversely, every catalogue record has a matching
  archived file or proof key.  Each appendix relation and every finite negative
  query used below references this complete catalogue.
\end{theorem}

\begin{proof}
  Enumerate the versioned spectrum and map manifests, compare their declared
  degree bounds with the $E_2$, map, and initial-$d_2$ files, and validate every
  foreign key in the propagated-output table.  Then compare normalized archive
  keys in both directions: no manifest or proof row may be omitted, and no
  stored record may lack an archive witness.  Finally check that the three
  manual inputs are absent from the machine-output side and occur exactly once
  in the separate manual ledger.
\end{proof}

% SOURCE: aimpaper/main.tex:2796-2800 and repeated rows at lines 2999, 3096, 3161.
\begin{proposition}[$x_{126,21}$ finite target ambiguity]
  \label{prop:x12621-ambiguity}
  \uses{def:near126-computation-input,def:appendix-evidence-record,thm:lin-computation-catalogues-complete}
  \notready
  For a near-$126$ computation input $I$, its known $d_4$ from $x_{126,21}$ has
  target candidate
  \[
    x_{125,25,2}+x_{125,25}+g^4\Delta h_1g
      +\epsilon d_0^2e_0gB_4,
    \qquad \epsilon\in\mathbb F_2,
  \]
  rather than a uniquely determined target.  All repeated table displays
  refer to one relation identifier with this two-element candidate set.
\end{proposition}

\begin{proof}
  Transcribe the prose and the three relevant table rows, normalize their class
  expressions, and verify that their fixed part agrees.  Encode the optional
  summand as the finite coefficient $\epsilon$, then link every display to the
  shared relation identifier.
\end{proof}

% SOURCE: aimpaper/main.tex:2738-2773, Table Table:Cnu126, PDF p.54.
\begin{proposition}[$S^0/\nu$ stem-$126$ table]
  \label{data:table-cnu126}
  \uses{def:near126-computation-input,def:appendix-evidence-record,thm:lin-computation-catalogues-complete}
  \notready
  For a near-$126$ computation input $I$, its Table 1 contains exactly 26
  nonempty rows in stem $126$ and filtrations
  $9\le s\le14$: 12 known incoming, 10 known outgoing, one permanent, and
  three unresolved rows, with no zero band.
\end{proposition}

\begin{proof}
  Parse lines 2738--2773, normalize the cell tags and expressions, and prove
  by finite enumeration that the four status counts sum to 26 and every
  printed class row occurs exactly once.
\end{proof}

% SOURCE: aimpaper/main.tex:2804-2848, Table Table:S122, PDF p.56.
\begin{proposition}[$S^0$ stem-$122$ table]
  \label{data:table-s122}
  \uses{def:near126-computation-input,def:appendix-evidence-record,thm:lin-computation-catalogues-complete}
  \notready
  For a near-$126$ computation input $I$, its Table 2 contains 35 nonempty
  rows: 16 incoming, 13 outgoing, five permanent, and one unresolved.  It also
  contains the zero bands $s=9$--$10$ and $s=0$--$7$.
\end{proposition}

\begin{proof}
  Transcribe lines 2804--2848 and enumerate the status constructors and two
  zero-band records.  Check the table's filtration coverage through $s=25$.
\end{proof}

% SOURCE: aimpaper/main.tex:2851-2903, Table Table:S123, PDF p.57.
\begin{proposition}[$S^0$ stem-$123$ table]
  \label{data:table-s123}
  \uses{def:near126-computation-input,def:appendix-evidence-record,thm:lin-computation-catalogues-complete}
  \notready
  For a near-$126$ computation input $I$, its Table 3 contains 42 nonempty
  rows: 22 incoming, 17 outgoing, two permanent, and one unresolved.  Its zero
  bands are $s=25$, $s=21$--$22$, and $s=0$--$7$.
\end{proposition}

\begin{proof}
  Transcribe lines 2851--2903, verify all row bidegrees and the three zero
  bands, and prove the status counts by finite enumeration.
\end{proof}

% SOURCE: aimpaper/main.tex:2908-2958, Table Table:S124.13, PDF p.58.
\begin{proposition}[$S^0$ stem-$124$ high-filtration table]
  \label{data:table-s124-high}
  \uses{def:near126-computation-input,def:appendix-evidence-record,thm:lin-computation-catalogues-complete}
  \notready
  For a near-$126$ computation input $I$, its Table 4 contains 43 nonempty rows
  for $13\le s\le25$: 24 incoming, 13 outgoing, and six permanent, with no
  unresolved row or zero band.
\end{proposition}

\begin{proof}
  Transcribe lines 2908--2958 and verify by enumeration that the three status
  counts sum to 43 and cover every printed filtration in the stated band.
\end{proof}

% SOURCE: aimpaper/main.tex:2962-2992, Table Table:S124.12, PDF p.59.
\begin{proposition}[$S^0$ stem-$124$ low-filtration table]
  \label{data:table-s124-low}
  \uses{def:near126-computation-input,def:appendix-evidence-record,thm:lin-computation-catalogues-complete}
  \notready
  For a near-$126$ computation input $I$, its Table 5 contains 23 nonempty rows
  for $s\le12$: 11 incoming, 11 outgoing, and one permanent.  It contains the
  zero band $s=0$--$5$.
\end{proposition}

\begin{proof}
  Transcribe lines 2962--2992, check the zero band and all class degrees, and
  prove the three status counts by finite enumeration.
\end{proof}

% SOURCE: aimpaper/main.tex:2994-3018, Table Table:S125.20, PDF p.59.
\begin{proposition}[$S^0$ stem-$125$ high-filtration table]
  \label{data:table-s125-high}
  \uses{def:near126-computation-input,def:appendix-evidence-record,thm:lin-computation-catalogues-complete,prop:x12621-ambiguity}
  \notready
  For a near-$126$ computation input $I$, its Table 6 contains 17 nonempty rows
  for $20\le s\le25$: six incoming, ten outgoing, and one permanent.  One
  outgoing row carries the finite ``possibly'' ambiguity and there is no zero
  band.
\end{proposition}

\begin{proof}
  Transcribe lines 2994--3018, connect the ambiguous row to the shared relation
  identifier, and enumerate the remaining statuses and bidegrees.
\end{proof}

% SOURCE: aimpaper/main.tex:3020-3078, Table Table:S125.19, PDF p.60.
\begin{proposition}[$S^0$ stem-$125$ low-filtration table]
  \label{data:table-s125-low}
  \uses{def:near126-computation-input,def:appendix-evidence-record,thm:lin-computation-catalogues-complete}
  \notready
  For a near-$126$ computation input $I$, its Table 7 contains 50 nonempty rows
  for $s\le19$: 18 incoming, 28 outgoing, two permanent, and two unresolved.
  It contains the zero band $s=0$--$4$.
\end{proposition}

\begin{proof}
  Transcribe lines 3020--3078 and prove the counts and filtration coverage.
  Preserve the zero band as a proposition-valued record because it is used in
  the main reduction to exclude low-weight $\lambda$-torsion.
\end{proof}

% SOURCE: aimpaper/main.tex:3081-3135, Table Table:S126.11, PDF p.61.
\begin{proposition}[$S^0$ stem-$126$ high-filtration table]
  \label{data:table-s126-high}
  \uses{def:near126-computation-input,def:appendix-evidence-record,thm:lin-computation-catalogues-complete,prop:x12621-ambiguity}
  \notready
  For a near-$126$ computation input $I$, its Table 8 contains 47 nonempty rows
  for $11\le s\le25$: 29 incoming, 15 outgoing, one permanent, and two
  unresolved.  One row has finite target ambiguity and there is no zero band.
\end{proposition}

\begin{proof}
  Transcribe lines 3081--3135, link the ambiguous display, and verify all four
  counts and the complete filtration interval by finite enumeration.
\end{proof}

% SOURCE: aimpaper/main.tex:3138-3171, Table Table:S126.10, PDF p.62.
\begin{proposition}[$S^0$ stem-$126$ low-filtration table]
  \label{data:table-s126-low}
  \uses{def:near126-computation-input,def:appendix-evidence-record,thm:lin-computation-catalogues-complete,prop:x12621-ambiguity}
  \notready
  For a near-$126$ computation input $I$, its Table 9 contains 26 nonempty rows
  for $s\le10$: 11 incoming, ten outgoing, one permanent, and four unresolved.
  One row has finite ambiguity; the zero band is $s=0$--$1$.  In particular the
  $h_6^2$ row is unresolved, not a known nonzero $d_7$.
\end{proposition}

\begin{proof}
  Transcribe lines 3138--3171, preserve each question-mark status, link the
  ambiguity relation, and verify counts and the single zero band.
\end{proof}

% SOURCE: aimpaper/main.tex:3173-3197, Table Table:S127.21, PDF p.62.
\begin{proposition}[$S^0$ stem-$127$ high-filtration table]
  \label{data:table-s127-high}
  \uses{def:near126-computation-input,def:appendix-evidence-record,thm:lin-computation-catalogues-complete}
  \notready
  For a near-$126$ computation input $I$, its Table 10 contains 17 nonempty rows
  for $21\le s\le25$: five incoming, 11 outgoing, and one permanent, with no
  unresolved row or zero band.
\end{proposition}

\begin{proof}
  Transcribe lines 3173--3197 and verify all class degrees, relation links, and
  status counts by finite enumeration.
\end{proof}

% SOURCE: aimpaper/main.tex:3200-3256, Table Table:S127.20, PDF p.63.
\begin{proposition}[$S^0$ stem-$127$ middle-filtration table]
  \label{data:table-s127-middle}
  \uses{def:near126-computation-input,def:appendix-evidence-record,thm:lin-computation-catalogues-complete}
  \notready
  For a near-$126$ computation input $I$, its Table 11 contains 50 nonempty rows
  for $10\le s\le20$: 20 incoming, 22 outgoing, and eight permanent, with no
  unresolved row or zero band.
\end{proposition}

\begin{proof}
  Transcribe lines 3200--3256 and prove by finite enumeration that all 50 rows
  have one of the three stated statuses and consistent relation identifiers.
\end{proof}

% SOURCE: aimpaper/main.tex:3258-3290, Table Table:S127.9, PDF p.64.
\begin{proposition}[$S^0$ stem-$127$ low-filtration table]
  \label{data:table-s127-low}
  \uses{def:near126-computation-input,def:appendix-evidence-record,thm:lin-computation-catalogues-complete}
  \notready
  For a near-$126$ computation input $I$, its Table 12 contains 25 nonempty rows
  for $s\le9$: four incoming, 14 outgoing, two permanent, and five unresolved.
  Its zero band is $s=0$.
\end{proposition}

\begin{proof}
  Transcribe lines 3258--3290, preserve all five unresolved statuses, and
  verify the counts, bidegrees, and zero band by finite enumeration.
\end{proof}

% SOURCE: aimpaper/main.tex:2738-3290, all tables; local PDF pp.54,56-64.
\begin{theorem}[All appendix rows are encoded]
  \label{thm:appendix-all-rows-encoded}
  \uses{def:near126-computation-input,thm:lin-computation-catalogues-complete,data:table-cnu126,data:table-s122,data:table-s123,data:table-s124-high,data:table-s124-low,data:table-s125-high,data:table-s125-low,data:table-s126-high,data:table-s126-low,data:table-s127-high,data:table-s127-middle,data:table-s127-low,prop:appendix-manual-inputs}
  \notready
  For a near-$126$ computation input $I$ satisfying the catalogue validation
  theorem, its twelve datasets contain exactly 401 nonempty class records and
  nine zero bands.  Among the nonempty records there are 178 incoming, 174 known
  outgoing, 31 permanent, and 18 unresolved statuses.  Every printed row and
  zero band appears exactly once, every duplicated differential is linked by
  one relation identifier, and no unprinted row is introduced.
\end{theorem}

\begin{proof}
  Concatenate the twelve finite tables and evaluate decidable checks for key
  uniqueness, source-line coverage, spectrum and bidegree consistency, status
  totals, and paired relation identifiers.  Compare the normalized
  serialization with the fixed source hashes and emit a coverage certificate
  listing every source row and zero band.
\end{proof}

% VERIFIED-IN: local source hashes recorded during audit.
\begin{proposition}[Primary-source snapshot]
  \label{prop:paper-source-snapshot}
  \uses{def:source-ref}
  \notready
  The audited snapshot records 66 PDF pages and SHA256 values
  \texttt{7cae269851a88d10dd194651dbf7497b75ef8b8b914bc1901dfa3734cd8096b4}
  for the PDF,
  \texttt{1125462bcae4a4ec56e3bfcaad15df4febf98757dfb83462b155af162c99c9e0}
  for \texttt{main.tex}, and
  \texttt{5eb83ef105ce9dbbda87a9f57726bb5c078308fd3718840ea29ee8100cbefabb}
  for \texttt{112.tex}.
\end{proposition}

\begin{proof}
  Recompute the three hashes, verify the PDF page count, and compare them with
  the stored source-reference fields before accepting any row-completeness
  certificate.
\end{proof}

\section{Near-126 finite evidence}

% SOURCE: aimpaper/main.tex:2127-2139, Remarks 7.4-7.5; Appendix Tables 3 and 5.
\begin{proposition}[$\theta_5$ order and torsion evidence]
  \label{evidence:theta5-order-torsion}
  \uses{def:near126-computation-input,def:synthetic-theta5-choice,thm:external-theta5-order-data,thm:external-synthetic-einfty-nu,def:external-evidence,data:table-s123,data:table-s124-low}
  \notready
  For a near-$126$ computation input $D$, the named Xu/IWX result and $D$'s
  finite Ext evidence supply: every classical
  $\theta_5$ has order two; the difference of two choices has filtration at
  least six; $\pi_{62,64}S^{0,0}$ has no $\lambda$-torsion; and
  $\pi_{125,130}S^{0,0}$ has no $\lambda$-torsion in the comparison range.
  These facts validate both the order-two hypothesis and the paper's
  one-$\lambda$ shift of the finite BJM/BX statement.
\end{proposition}

\begin{proof}
  Load the located homotopy-group structure theorem and order computation, then
  evaluate the synthetic $E_\infty$ formula on the complete table data in the
  two stated bidegrees.  Exhaustively exclude a class with the required
  torsion length and record the search range in the evidence certificate.
\end{proof}

% SOURCE: aimpaper/main.tex:2156-2177, Fact fact:theta5sqAF and following remark; Appendix Tables 5,6,7,9.
\begin{proposition}[Core near-$126$ evidence]
  \label{evidence:near126-core}
  \uses{def:near126-computation-input,def:near126-conditions,data:table-s124-low,data:table-s125-high,data:table-s125-low,data:table-s126-low,thm:appendix-all-rows-encoded,prop:lin-computation-provenance}
  \notready
  For a near-$126$ computation input $D$, its explicit computation evidence
  supplies:
  \begin{enumerate}
    \item $x_{126,8,4}+x_{126,8}$ survives through $E_6$;
    \item $h_1h_4x_{109,12}$ supports no outgoing differential, and the
      exhaustive search leaves exactly two possible incoming killers:
      $d_6(x_{126,8,4}+x_{126,8})$ and $d_{12}(h_6^2)$;
    \item $h_0^2x_{124,8}$ is permanent;
    \item $g^4\Delta h_1g$ is the unique class in its bidegree surviving to
      $E_5$;
    \item $d_3(x_{126,6})$ is one of the two printed nonzero target sums and
      therefore cannot kill $h_1h_4x_{109,12}$.
  \end{enumerate}
  Item (2) deliberately does not call the class permanent before both possible
  incoming differentials are excluded.
\end{proposition}

\begin{proof}
  Run typed queries over the complete four table ranges and their differential
  relation graph.  Check survival, outgoing and incoming ranges separately,
  preserve the finite ambiguity of $d_3(x_{126,6})$, and attach the Zenodo
  disproof identifiers for every excluded potential differential.
\end{proof}

% SOURCE: aimpaper/main.tex:2221-2258, proofs of Lemmas lem:equistate4 and lem:equistate5.
\begin{proposition}[Representative-indeterminacy evidence]
  \label{evidence:near126-indeterminacy}
  \uses{def:near126-computation-input,evidence:theta5-order-torsion,data:table-s124-high,data:table-s124-low,thm:appendix-all-rows-encoded}
  \notready
  For a near-$126$ computation input $D$, its complete stem-$124$ data supply
  that filtration-11 and -12 cycles relevant
  to $C_5$ are killed by $d_2$ or $d_3$, filtration-13 correction terms are
  annihilated by $h_1$, and all remaining corrections begin in filtration at
  least 15 after multiplication by $\lambda^3\eta$.
\end{proposition}

\begin{proof}
  Enumerate every cycle in stem 124 and filtrations 11--13 from both table
  bands, join its outgoing relation and product-by-$h_1$ record, and verify the
  stated alternative for each.  Completeness excludes an unlisted correction.
\end{proof}

% SOURCE: aimpaper/main.tex:2323-2354, exhaustive theta-square and eta-square analysis; tmf citation.
\begin{proposition}[$\theta_5^2$ and $tmf$ exhaustive evidence]
  \label{evidence:theta5-square-tmf}
  \uses{def:near126-computation-input,evidence:near126-core,evidence:near126-indeterminacy,data:table-s125-low,data:table-s125-high,def:external-result,def:external-evidence}
  \notready
  For a near-$126$ computation input $D$, its explicit table and $tmf$ inputs
  supply:
  \begin{enumerate}
    \item the only possibilities for $\theta_5^2$ are
      $\lambda^6[h_0^2x_{124,8}]$,
      $\lambda^9[e_0\Delta h_6g]$, or a $\lambda^{10}$-multiple;
    \item $h_1e_0\Delta h_6g=0$;
    \item the only relevant high-filtration $\lambda^{10}$-multiple is the
      class represented by $\lambda^{20}g^4\Delta h_1g$;
    \item its $tmf$ Hurewicz image is nonzero, while every $\theta_5$ maps to
      zero in $tmf$;
    \item if $C_3$ fails, the unique alternative incoming differential is
      $d_6(x_{126,8,4}+x_{126,8})=h_1h_4x_{109,12}$.
  \end{enumerate}
\end{proposition}

\begin{proof}
  Query the complete filtration ranges and product table to produce the three
  alternatives and unique surviving high class.  Attach independent, located
  $tmf$ Hurewicz evidence for its image and the image of $\theta_5$.  Finally
  use the incoming-differential completeness certificate for the stated
  alternative.  Each negative conclusion retains its finite search bounds.
\end{proof}

% SOURCE: aimpaper/main.tex:2377-2385, Fact fact:x1239; Tables S123 and S125.19.
\begin{proposition}[$x_{123,9}$ computation evidence]
  \label{evidence:x1239}
  \uses{def:near126-computation-input,data:table-s123,data:table-s125-low,thm:appendix-all-rows-encoded}
  \notready
  For a near-$126$ computation input $D$, the class
  $x_{123,9}+h_0x_{123,8}$ survives through $E_{12}$ and has no incoming
  differential in $D$, and
  \[
    d_2(x_{125,8})=
      h_1(x_{123,9}+h_0x_{123,8})+h_0^2x_{124,8}.
  \]
\end{proposition}

\begin{proof}
  Retrieve the named rows and shared differential relation, then use complete
  source ranges to certify the negative incoming claim through the relevant
  maximum length.
\end{proof}

% SOURCE: aimpaper/main.tex:2466-2468, Fact fact:h02x1259; Table S125.19.
\begin{proposition}[$h_0^2x_{125,9,2}$ evidence]
  \label{evidence:h02x1259}
  \uses{def:near126-computation-input,data:table-s125-low,thm:appendix-all-rows-encoded}
  \notready
  For a near-$126$ computation input $D$, the class $h_0^2x_{125,9,2}$
  survives through $E_5$ and no incoming
  differential hits it.  Its printed $d_5?$ status remains unresolved and is
  not asserted nonzero.
\end{proposition}

\begin{proof}
  Read the typed row and query every possible incoming source in the complete
  adjacent-stem ranges.  Preserve the unresolved-outgoing constructor for the
  $d_5$ column.
\end{proof}

% SOURCE: aimpaper/main.tex:2613-2615, Fact fact:h1x1217; Table S122.
\begin{proposition}[$h_1x_{121,7}$ evidence]
  \label{evidence:h1x1217}
  \uses{def:near126-computation-input,data:table-s122,thm:appendix-all-rows-encoded}
  \notready
  For a near-$126$ computation input $D$, the class $h_1x_{121,7}$ survives
  through $E_6$ and has no incoming
  differential in the complete range.
\end{proposition}

\begin{proof}
  Query its table row and all degree-compatible incoming relation identifiers;
  use table completeness to certify the negative result.
\end{proof}

% SOURCE: aimpaper/main.tex:2617-2673, missing obstruction check in Lemma lem:nuext125.
\begin{proposition}[Obstruction vanishing for the mod-$\lambda^5$ lift]
  \label{evidence:nu-extension-obstruction-vanishing}
  \uses{def:near126-computation-input,data:table-s125-low,data:table-s126-low,data:table-cnu126,thm:appendix-all-rows-encoded}
  \notready
  For a near-$126$ computation input $D$ and the relation produced modulo
  $\lambda^3$ by the $S^0/\nu_{\mathrm{Hopf}}$
  $d_3$, every filtration-10 correction that could obstruct a lift modulo
  $\lambda^5$ either has no detected homotopy class or can be killed by changing
  the lift.  Therefore a lift exists with no additional $[\lambda x]$ term.
\end{proposition}

\begin{proof}
  Enumerate the complete filtration-10 source group, compute the relevant
  $\rho$ obstruction and each candidate's synthetic survival from the table
  relations, and give an explicit correction or vanishing certificate for
  every element.  This evidence is independent of the formal properties of
  $\rho$; it fills the unsupported ``upon inspection'' step at
  \texttt{aimpaper/main.tex:2661--2665}.
\end{proof}

% SOURCE: aimpaper/main.tex:2689-2696, Fact fact:stem122; Table S122.
\begin{proposition}[Stem-$122$ permanent-class evidence]
  \label{evidence:stem122-permanent}
  \uses{def:near126-computation-input,data:table-s122,thm:appendix-all-rows-encoded}
  \notready
  For a near-$126$ computation input $D$, the classes $h_6Md_0$ and
  $h_5x_{91,11}$ are permanent in the classical Adams spectral sequence
  encoded by $D$.
\end{proposition}

\begin{proof}
  Retrieve their rows, certify no outgoing or incoming differential through
  the complete source ranges, and validate the permanent status against the
  program artifact identifiers.
\end{proof}

% SOURCE: aimpaper/main.tex:2707-2728, two ``upon inspection'' steps in the proof of Proposition prop:state5false.
\begin{proposition}[Stem-$122$ candidate and product exhaustion]
  \label{evidence:stem122-product-exhaustion}
  \uses{def:near126-computation-input,data:table-s122,data:table-s125-low,evidence:stem122-permanent,thm:appendix-all-rows-encoded,def:external-evidence}
  \notready
  For a near-$126$ computation input $D$, in the bidegree and
  $\lambda$-quotient occurring in the proof of
  Proposition \texttt{prop:state5false}, every nonzero representative of
  $[\lambda^4h_1x_{121,7}][h_0]$ of Adams filtration at most $12$ is exactly
  one of
  \[
    \lambda^6[h_6Md_0]\quad\text{or}\quad
    \lambda^7[h_5x_{91,11}].
  \]
  The typed product and differential evidence further shows, in either case,
  that $\lambda^4[h_1h_4x_{109,12}]$ is a
  $\lambda[h_2]$-multiple in the untruncated synthetic sphere.  This record
  includes the filtration-$13$ candidate search and the $d_2/d_4$ exclusions
  used in lines 2716--2728.
\end{proposition}

\begin{proof}
  Enumerate the complete stem-$122$ permanent representatives through
  filtration $12$, evaluate the typed $h_0$ and $h_2$ products, and transport
  each known differential through the shared relation identifiers.  Query the
  complete stem-$125$ ranges for filtration-$13$ corrections and prove that
  each is killed by the recorded $d_2$ or $d_4$.  The resulting finite
  certificate has exactly the two source cases and the asserted common
  divisibility conclusion.
\end{proof}

% SOURCE: aimpaper/main.tex:1398-1405, Example exam:extonEn(2), for the
% identity \widehat{\nu_{\mathrm{Hopf}}}=[h_2] and its grading;
% aimpaper/main.tex:929-961, Notation not:fhat, for normalization.
\begin{proposition}[Normalized Hopf-map detection]
  \label{evidence:normalized-hopf-detection}
  \uses{def:synthetic-hat-map,def:adams-detection,thm:external-synthetic-triangle-lift,def:cofiber-triangle,def:external-evidence}
  \notready
  Explicit detection evidence retains one chosen application of the synthetic
  triangle-lift theorem to the appropriate rotation of the classical Hopf
  cofiber triangle.  It identifies that stored triangle's Hopf component with
  \[
    \widehat{\nu_{\mathrm{Hopf}}}=[h_2]:S^{3,4}\longrightarrow S^{0,0},
  \]
  where $[h_2]$ is detected by $h_2$ in the source's stated synthetic
  tridegree.  With this fixed normalization,
  $\nu(\nu_{\mathrm{Hopf}})=\lambda[h_2]$.
\end{proposition}

\begin{proof}
  Use the explicit normalization and tridegree calculation in Example
  \texttt{exam:extonEn(2)} at \texttt{aimpaper/main.tex:1398--1405}.  Store
  the entire chosen lifted-triangle witness, the detected class, the equality
  with its Hopf component, and the factorization by $\lambda$ as one
  source-bearing evidence value.
\end{proof}

% SOURCE: aimpaper/main.tex:2698-2778 and Table Table:Cnu126.
\begin{proposition}[$S^0/\nu$ short-incoming exclusion]
  \label{evidence:cnu126-short-incoming-exclusion}
  \uses{def:near126-computation-input,data:table-cnu126,thm:appendix-all-rows-encoded,prop:lin-computation-provenance}
  \notready
  For a near-$126$ computation input $D$, in its classical Adams spectral
  sequence of $S^0/\nu_{\mathrm{Hopf}}$, the class
  $h_1h_4x_{109,12}[0]$ is not hit by any $d_r$ with $r\le5$.  This negative
  statement is a range-completeness query over all 26 table rows, not a claim
  read from a single row.
\end{proposition}

\begin{proof}
  Enumerate every degree-compatible source in filtrations 9--14, follow its
  relation status, and prove none has target the named bottom-cell class with
  length at most five.  Table and artifact completeness make the finite
  negative query exhaustive.
\end{proof}
